% Kurzexposé -- genau eine Seite, Anlage zur Mail an die Betreuung.
% Bauen: ./build.sh expose   ->  expose.pdf
%
% Bewusst eigenständig: kein Literaturverzeichnis, keine Kapitelnummern.
% Inhaltlich auf dem Stand von report/chapters/03-05.

\documentclass[10pt,a4paper,ngerman]{article}

\usepackage{iftex}
\usepackage[T1]{fontenc}
\ifPDFTeX
  \usepackage[utf8]{inputenc}
\fi
\usepackage{lmodern}
\usepackage{microtype}

% XeTeX mit T1-Schriften setzt das direkte Zeichen ß als "SS".
% ��ber newunicodechar wird es auf \ss{} umgebogen -- wirkt auch unter pdfTeX.
\usepackage{newunicodechar}
\newunicodechar{ß}{\ss{}}
\usepackage[ngerman]{babel}

\usepackage[a4paper,top=1.7cm,bottom=1.6cm,left=1.9cm,right=1.9cm]{geometry}
\usepackage{booktabs}
\usepackage{array}
\usepackage{enumitem}
\usepackage{xcolor}

% Titelblatt-Daten (Matrikelnummer) stehen lokal in private.tex und werden nie
% versioniert. Fehlt die Datei, bleibt der Platzhalter stehen.
\IfFileExists{private.tex}
  {\input{private.tex}}
  {\newcommand{\matrikelnummer}{[Matrikelnummer]}}

\setlength{\parindent}{0pt}
\setlength{\parskip}{2.5pt}
\pagestyle{empty}
\setlist[itemize]{leftmargin=1em,itemsep=0.3pt,topsep=1pt,parsep=0pt}
\renewcommand{\arraystretch}{1.0}

% Abschnittskopf: klein, fett, dünne Linie darunter
\newcommand{\abschnitt}[1]{%
  \vspace{3.5pt}%
  {\small\bfseries #1}\par
  \vspace{-3pt}{\color{black!30}\rule{\linewidth}{0.4pt}}\par
  \vspace{-1pt}}

\begin{document}

{\footnotesize\bfseries Kurzexposé\par}
\vspace{1pt}
{\Large\bfseries Konsistenz über Quellen\par}
\vspace{0.5pt}
{\normalsize Abweichungen zwischen Studieninformationen automatisch finden und bewerten\par}
\vspace{2pt}
{\footnotesize
The Trung Do \textbullet{} M.Sc. Data Analytics, Matrikelnummer \matrikelnummer{} \textbullet{}
Projekt 29 (Team Gamma) \textbullet{} Betreuung: Dr. Sebastian Busse\\
Stand: 21. November 2026 \textbullet{} Arbeitstitel \textbullet{} Prüfungsleistung: Bericht
10--15 Seiten mit Programmcode, ohne Vortrag\par}

% ---------------------------------------------------------------------
\abschnitt{Fragestellung}

Dieselben Sachverhalte -- Fristen, Leistungspunkte, Modulnummern, Ansprechpartner --
stehen an mehreren Stellen: in der Speziellen Ordnung, in den Allgemeinen Bestimmungen,
auf der Studiengangsseite, im Studienführer und auf der Fachbereichsseite. Ob sie
übereinstimmen, ist offen.

\textbf{Leitfrage.} Welche Sachverhalte werden in mehreren Quellen unterschiedlich
dargestellt, und wie zuverlässig lässt sich das automatisch finden?

\textbf{H1 (inhaltlich).} Ein großer Teil der Abweichungen ist kein echter Widerspruch,
sondern Folge unterschiedlicher Fassungen oder Zielgruppen; echte Widersprüche sind
seltener, treten aber systematisch bei Fristen auf.

\textbf{H2 (methodisch).} Sprachmodellgestützte Extraktion erreicht hohe Precision bei
begrenztem Recall; regelbasierte Normalisierung erhöht den Recall, ohne die Precision zu
senken.

% ---------------------------------------------------------------------
\abschnitt{Daten}

Öffentliche Quellen, \textbf{ein} Stichtag, ein Hash je Dokument, ein Datenblatt mit
Umfang und Lücken.

\begin{itemize}
  \item Spezielle Ordnung M.Sc. Data Analytics, Fassung des 3. Änderungsbeschlusses
        (gültig ab SoSe 2026) \emph{und} die drei früheren Fassungen -- erst sie trennen
        Versionsartefakte von echten Widersprüchen.
  \item Allgemeine Bestimmungen (3. Änderungsfassung): dieselbe Frist steht dort allgemein
        und in der Speziellen Ordnung speziell.
  \item Studiengangsseite, Studienführer der Zentralen Studienberatung, Fachbereichsseite;
        für datierbare Fassungen die Wayback Machine (CDX-API).
\end{itemize}

Nicht erfasst: Login-Bereiche, Scans ohne Textebene, dynamische Inhalte. Korpus und
Faktkatalog (rund zwölf Typen) lege ich mit Team Gamma bis Mitte Dezember fest.

% ---------------------------------------------------------------------
\abschnitt{Methode}

\textbf{Analyseeinheit.} Sachverhalt als Tupel aus Entität, Attribut, Wert,
Gültigkeitsintervall, Quelle und Fundstelle -- die Definition, die
\glqq{}derselbe Sachverhalt\grqq{} überhaupt prüfbar macht.

\textbf{Extraktion mit Zitatpflicht.} Sprachmodell mit festem Schema; jeder Wert muss eine
Fundstelle mitliefern, Werte ohne Beleg werden verworfen.

\textbf{Normalisierung.} Datum (ISO 8601), Zahl, Einheit, Entitätsalias, Zielgruppe --
sonst entstehen falsche Widersprüche.

\textbf{Typologie.} Echter Widerspruch, Versionsartefakt, Scope-Unterschied, Formulierung,
Extraktionsfehler.

\textbf{Evaluation.} Stratifizierte Stichprobe, Goldstandard mit zwei Durchgängen,
Krippendorffs $\alpha$; Precision, Recall, F1 je Sachverhaltstyp mit
Wilson-Konfidenzintervall. Die Übereinstimmung Modell/Mensch wird vor dem Lauf über den
vollen Korpus berichtet; Fehleranalyse über 20--30 Fälle.

% ---------------------------------------------------------------------
\abschnitt{Zeitplan}

{\small
\begin{tabular}{@{}p{3.1cm}p{3.4cm}p{8.6cm}@{}}
\toprule
Phase & Zeitraum & Ziel \\
\midrule
P1 Grundlage     & 01.12.--20.12.2026 & Literatur, Hypothesen, Korpus, Datenblatt v1 \\
P2 Plan/Prototyp & 04.01.--17.01.2027 & Projektplan, Meilensteingespräch, Prototyp (3 Dokumente) \\
P3 Pipeline      & 18.01.--14.02.2027 & Extraktion, Normalisierung, Vergleich, Goldstandard \\
P4 Evaluation    & 15.02.--14.03.2027 & Annotation, Metriken, Fehleranalyse \\
P5 Schreiben     & 15.03.--28.03.2027 & Bericht vollständig \\
P6 Abgabe        & 29.03.--01.04.2027 & Endfassung, Code, Abgabe \\
\bottomrule
\end{tabular}}

\vspace{1pt}
{\footnotesize 21.12.2026--03.01.2027: Weihnachtspause, Korpus eingefroren.
\textbf{Bitte:} kurze Rückmeldung, ob Fragestellung, Datenzuschnitt und Evaluationsplan so
tragen -- bevor ich Kapitel 5 und 6 schreibe.\par}

\end{document}
